
\documentclass{article}
\usepackage{colortbl}
\usepackage{makecell}
\usepackage{multirow}
\usepackage{supertabular}

\begin{document}

\newcounter{utterance}

\centering \large Interaction Transcript for game `cladder', experiment `full\_v1.5\_default', episode 3456 with selfrefine{-}Qwen3.5{-}27B.
\vspace{24pt}

{ \footnotesize  \setcounter{utterance}{1}
\setlength{\tabcolsep}{0pt}
\begin{supertabular}{c@{$\;$}|p{.15\linewidth}@{}p{.15\linewidth}p{.15\linewidth}p{.15\linewidth}p{.15\linewidth}p{.15\linewidth}}
    \# & \multicolumn{2}{c}{Player} && \multicolumn{2}{c}{Game Master} \\
    \hline

    \theutterance \stepcounter{utterance}  
    & & & \multicolumn{4}{p{0.6\linewidth}}{
        \cellcolor[rgb]{0.9,0.9,0.9}{
            \makecell[{{p{\linewidth}}}]{
                \texttt{\tiny{[P1$\langle$GM]}}
                \texttt{You are an expert in causal inference. The following question is not a typical commonsense query, but rather a meticulously designed question created by a professor specializing in causal inference, intended to assess the students' mastery of the course content.} \\
\texttt{Imagine a self{-}contained, hypothetical world with only the following conditions, and without any unmentioned factors or causal relationships: Solar eclipse has a direct effect on Alice arriving to school. Traffic has a direct effect on Alice arriving to school. Method 1: We look at how solar eclipse correlates with Alice arriving to school case by case according to traffic. Method 2: We look directly at how solar eclipse correlates with Alice arriving to school in general. To understand how solar eclipse affects Alice arriving to school, is it more correct to use the Method 1 than Method 2?} \\
\texttt{Start your answer with 'Yes' or 'No', followed by additional reasoning or evidence to support your explanation.} \\
            }
        }
    }
     \\ \\

    \theutterance \stepcounter{utterance}  
    & & \multicolumn{4}{p{0.6\linewidth}}{
        \cellcolor[rgb]{0.95,0.95,0.95}{
            \makecell[{{p{\linewidth}}}]{
                \texttt{\tiny{[P1$\leftrightarrow$P]}}
                \texttt{{[}DRAFT{]}} \\
\texttt{No} \\
\\ 
\texttt{In the hypothetical world described, the causal structure is defined such that "Solar eclipse" (\$X\$) and "Traffic" (\$Z\$) are both direct causes of "Alice arriving to school" (\$Y\$), with no causal link between \$X\$ and \$Z\$ (since no unmentioned factors or relationships exist). This structure represents a **fork** or a **common effect** scenario where \$X\$ and \$Z\$ are independent causes of \$Y\$.} \\
\\ 
\texttt{Method 2 (looking at the general correlation) estimates the **total effect** of the solar eclipse on Alice's arrival. Since there is no confounding variable (a common cause of both the eclipse and traffic) and no mediator (traffic is not caused by the eclipse), the marginal association between the solar eclipse and Alice's arrival accurately reflects the direct causal effect.} \\
\\ 
\texttt{Method 1 (stratifying by traffic) attempts to control for traffic. However, because traffic is not a confounder (it does not cause the solar eclipse) and is not a mediator (the solar eclipse does not cause traffic), conditioning on it is unnecessary. In fact, if there were any unmeasured common causes of traffic and arrival (which are explicitly excluded here), or if traffic were a collider (which it is not in this specific description, as it is a direct cause, not an effect of both), conditioning could introduce bias. More importantly, in this specific "no unmentioned factors" setup, stratifying by an independent cause (\$Z\$) simply partitions the data without changing the causal estimate, but it reduces statistical efficiency (increases variance) compared to the marginal estimate. Therefore, Method 2 is the more direct and statistically efficient approach to understanding the effect, making Method 1 not "more correct" but rather redundant or less efficient. The question asks if Method 1 is *more* correct; it is not.} \\
            }
        }
    }
    & \\ \\

    \theutterance \stepcounter{utterance}  
    & & \multicolumn{4}{p{0.6\linewidth}}{
        \cellcolor[rgb]{0.95,0.95,0.95}{
            \makecell[{{p{\linewidth}}}]{
                \texttt{\tiny{[P1$\leftrightarrow$P]}}
                \texttt{{[}SELF{-}FEEDBACK{]}} \\
\texttt{FORMAT: ok} \\
\texttt{RULES: ok} \\
\texttt{STRATEGY: ok} \\
\texttt{REASON: The answer correctly starts with "No" as required, accurately identifies the causal structure (independent causes), and logically explains that stratifying by a non{-}confounder (Method 1) is less efficient than the marginal estimate (Method 2) in this specific setup.} \\
            }
        }
    }
    & \\ \\

    \theutterance \stepcounter{utterance}  
    & \multicolumn{4}{p{0.6\linewidth}}{
        \cellcolor[rgb]{0.561,0.737,0.561}{
            \makecell[{{p{\linewidth}}}]{
                \texttt{\tiny{[P1$\rangle$GM]}}
                \texttt{No} \\
\\ 
\texttt{In the hypothetical world described, the causal structure consists of "Solar eclipse" (\$X\$) and "Traffic" (\$Z\$) as independent direct causes of "Alice arriving to school" (\$Y\$), with no causal link between \$X\$ and \$Z\$ and no unmentioned confounders. This configuration represents two independent causes acting on a common effect.} \\
\\ 
\texttt{Method 2 (looking at the general correlation) estimates the total causal effect of the solar eclipse on Alice's arrival. Since there is no confounding (no common cause of \$X\$ and \$Y\$) and no mediation (traffic is not on the causal path from eclipse to arrival), the marginal association directly reflects the causal effect.} \\
\\ 
\texttt{Method 1 (stratifying by traffic) conditions on \$Z\$. Because \$Z\$ is not a confounder, conditioning on it does not remove any bias. Instead, it merely partitions the data, which reduces the sample size in each stratum and increases the variance of the estimate without improving accuracy. Therefore, Method 1 is statistically less efficient than Method 2. Since Method 1 offers no advantage in correctness and sacrifices efficiency, it is not "more correct" to use Method 1 than Method 2.} \\
            }
        }
    }
    & & \\ \\

    \theutterance \stepcounter{utterance}  
    & & & \multicolumn{2}{p{0.3\linewidth}}{
        \cellcolor[rgb]{0.9,0.9,0.9}{
            \makecell[{{p{\linewidth}}}]{
                \texttt{\tiny{[GM$|$GM]}}
                \texttt{No} \\
            }
        }
    }
    & & \\ \\

    \theutterance \stepcounter{utterance}  
    & & & \multicolumn{2}{p{0.3\linewidth}}{
        \cellcolor[rgb]{0.9,0.9,0.9}{
            \makecell[{{p{\linewidth}}}]{
                \texttt{\tiny{[GM$|$GM]}}
                \texttt{no} \\
            }
        }
    }
    & & \\ \\

    \theutterance \stepcounter{utterance}  
    & & & \multicolumn{2}{p{0.3\linewidth}}{
        \cellcolor[rgb]{0.9,0.9,0.9}{
            \makecell[{{p{\linewidth}}}]{
                \texttt{\tiny{[GM$|$GM]}}
                \texttt{game\_result = WIN} \\
            }
        }
    }
    & & \\ \\

\end{supertabular}
}

\end{document}
